\documentclass[11pt,a4paper]{ctexart}
\usepackage[a4paper,margin=1.68cm]{geometry}
\usepackage{amsmath,amssymb,mathtools}
\usepackage{booktabs,tabularx,longtable,array}
\usepackage{enumitem}
\usepackage{tikz}
\usetikzlibrary{arrows.meta,positioning,shapes.geometric,fit,calc,matrix,backgrounds}
\usepackage[most]{tcolorbox}
\usepackage{xcolor}
\usepackage{fancyhdr}
\usepackage{hyperref}
\usepackage{microtype}
\usepackage{pifont}
\usepackage{caption}
\newcolumntype{Y}{>{\raggedright\arraybackslash}X}
\setlength{\parindent}{2em}
\setcounter{secnumdepth}{0}
\setlength{\parskip}{0.28em}
\setlist[itemize]{itemsep=0.14em,topsep=0.24em,leftmargin=2em}
\setlist[enumerate]{itemsep=0.14em,topsep=0.24em,leftmargin=2em}
\hypersetup{colorlinks=true,linkcolor=blue!45!black,urlcolor=blue!50!black,pdftitle={英语语言学习与考试统一方法论}}

\definecolor{mainblue}{RGB}{43,85,135}
\definecolor{deepblue}{RGB}{28,60,98}
\definecolor{softblue}{RGB}{238,245,252}
\definecolor{softgreen}{RGB}{238,249,243}
\definecolor{softorange}{RGB}{253,246,233}
\definecolor{softred}{RGB}{253,239,239}
\definecolor{softgray}{RGB}{246,247,249}
\definecolor{deepgreen}{RGB}{35,105,75}
\definecolor{deeporange}{RGB}{165,98,22}
\definecolor{deepred}{RGB}{150,55,55}
\definecolor{purple}{RGB}{94,66,140}
\definecolor{softpurple}{RGB}{245,241,251}
\definecolor{teal}{RGB}{38,112,115}
\definecolor{softteal}{RGB}{238,248,248}
\definecolor{gold}{RGB}{156,120,28}
\definecolor{softgold}{RGB}{252,249,235}

\pagestyle{fancy}
\fancyhf{}
\lhead{英语语言学习与考试统一方法论}
\rhead{从语言系统到任务执行}
\cfoot{\thepage}
\renewcommand{\headrulewidth}{0.4pt}
\setlength{\headheight}{14pt}

\newtcolorbox{corebox}[1][]{colback=softblue,colframe=mainblue,title=#1,fonttitle=\bfseries,boxrule=0.8pt,arc=2mm}
\newtcolorbox{exam}[1][]{colback=softorange,colframe=deeporange,title=#1,fonttitle=\bfseries,boxrule=0.8pt,arc=2mm}
\newtcolorbox{warn}[1][]{colback=softred,colframe=deepred,title=#1,fonttitle=\bfseries,boxrule=0.8pt,arc=2mm}
\newtcolorbox{eng}[1][]{colback=softgreen,colframe=deepgreen,title=#1,fonttitle=\bfseries,boxrule=0.8pt,arc=2mm}
\newtcolorbox{method}[1][]{colback=softgray,colframe=black!45,title=#1,fonttitle=\bfseries,boxrule=0.6pt,arc=2mm}
\newtcolorbox{bridge}[1][]{colback=softpurple,colframe=purple,title=#1,fonttitle=\bfseries,boxrule=0.8pt,arc=2mm}
\newtcolorbox{statebox}[1][]{colback=softteal,colframe=teal,title=#1,fonttitle=\bfseries,boxrule=0.8pt,arc=2mm}
\newtcolorbox{history}[1][]{colback=softgold,colframe=gold,title=#1,fonttitle=\bfseries,boxrule=0.8pt,arc=2mm}

\newcommand{\kw}[1]{\textbf{#1}}
\newcommand{\term}[2]{\textbf{#1（#2）}}
\newcommand{\yes}{\textcolor{deepgreen}{\ding{51}}}
\newcommand{\no}{\textcolor{deepred}{\ding{55}}}

% ---- Figure grammar: direct topology & semantic color system ----
\tikzset{
  cnode/.style={draw=mainblue,rounded corners=1.8mm,minimum width=2.65cm,minimum height=0.92cm,
    align=center,inner sep=3pt,font=\small,fill=softblue,line width=0.8pt},
  cnodeblue/.style={cnode,draw=mainblue,fill=softblue},
  cnodeg/.style={cnode,draw=deepgreen,fill=softgreen,font=\small\bfseries},
  cnodegreen/.style={cnode,draw=deepgreen,fill=softgreen,font=\small\bfseries},
  cnodeo/.style={cnode,draw=deeporange,fill=softorange},
  cnodeorange/.style={cnode,draw=deeporange,fill=softorange},
  cnodep/.style={cnode,draw=purple,fill=softpurple},
  cnodepurple/.style={cnode,draw=purple,fill=softpurple},
  cnodet/.style={cnode,draw=teal,fill=softteal},
  cnodeteal/.style={cnode,draw=teal,fill=softteal},
  cnodered/.style={cnode,draw=deepred,fill=softred,font=\small\bfseries},
  cnodegray/.style={cnode,draw=black!45,fill=softgray},
  mainflow/.style={-{Latex[length=2.2mm,width=1.5mm]},line width=1.0pt,draw=deepblue},
  altflow/.style={-{Latex[length=2.0mm,width=1.4mm]},line width=0.9pt,draw=deepgreen},
  returnflow/.style={-{Latex[length=2.0mm,width=1.4mm]},densely dashed,line width=0.85pt,draw=purple},
  dangerflow/.style={-{Latex[length=2.0mm,width=1.4mm]},line width=0.9pt,draw=deepred},
  optflow/.style={-{Latex[length=1.8mm,width=1.3mm]},dotted,line width=0.85pt,draw=teal},
  labelnote/.style={font=\scriptsize\bfseries,fill=white,inner sep=1.5pt,rounded corners=0.5mm,text=black!80,draw=black!25,line width=0.4pt},
  boxgroup/.style={draw=black!25,dashed,rounded corners=2mm,fill=black!2,inner sep=6pt}
}

\title{\vspace{-1.15cm}\textbf{英语语言学习与考试统一方法论}\\[0.35em]
\large 从语言系统、在线处理到考试任务执行的统一心智模型}
\author{英语学科 Canonical Subject Map}
\date{}

\begin{document}
\maketitle

\begin{corebox}[总纲]
英语学习的目标不是积累若干相互分离的“听力技巧、阅读技巧、写作模板”，而是建立一套能够在不同任务中稳定运行的语言系统。
\begin{center}
\begin{tikzpicture}[node distance=0.40cm]
\node[cnodeblue,minimum width=2.6cm] (m1) {1. 语言资源\\[-1pt]\scriptsize 形式/词法/语义/语篇};
\node[cnodeorange,right=of m1,minimum width=2.6cm] (m2) {2. 在线处理\\[-1pt]\scriptsize 感知/解析/提取/监控};
\node[cnodeteal,right=of m2,minimum width=2.6cm] (m3) {3. 沟通任务\\[-1pt]\scriptsize 接收/产出/互动/中介};
\node[cnodegreen,right=of m3,minimum width=2.8cm] (m4) {4. 监控与修复\\[-1pt]\scriptsize 风险发现与策略控制};

\draw[mainflow] (m1) -- (m2);
\draw[mainflow] (m2) -- (m3);
\draw[altflow] (m3) -- (m4);
\end{tikzpicture}
\end{center}
考研英语、IELTS、TOEFL，以及以 CEFR A1--C2 为参照的各类课程与考试，都只是对这套系统的不同取样。题型可以变化，底层能力结构不随题型变化。
\end{corebox}

\begin{warn}[本册的边界]
本册不提供某一考试的题型技巧，不重复完形、阅读、新题型、写作、听力或口语专题中的细节方法。它只拥有三类稳定内容：
\begin{enumerate}
  \item 英语能力由哪些可区分的部件组成；
  \item 不同语言任务怎样调用这些部件；
  \item 学习者如何定位断点、选择训练目标，并把专项能力重新带回完整任务。
\end{enumerate}
\end{warn}

\tableofcontents
\newpage

\section{第 0 章：为什么需要一张跨考试的英语总图}

\subsection{0.1 “考研英语”“雅思英语”“托福英语”不是三种英语}
同一个学习者可以面对完全不同的考试形式：考研英语偏重阅读、翻译与书面产出；IELTS 同时直接测量听、说、读、写；TOEFL 更强调学术场景中的接收、产出和综合任务。考试形式不同，但都必须调用同一套英语资源：词汇、语法、语篇、语用，以及在时间压力下识别、理解、计划和输出语言的能力。

因此，考试名称不应该成为知识体系的最高层。更合理的结构是：
\begin{center}
\begin{tikzpicture}[node distance=0.40cm]
\node[cnodeblue,minimum width=2.6cm] (a1) {1. 语言系统\\[-1pt]\scriptsize 资源与在线处理};
\node[cnodeorange,right=of a1,minimum width=2.6cm] (a2) {2. 语言活动\\[-1pt]\scriptsize 接收/产出/互动/中介};
\node[cnodeteal,right=of a2,minimum width=2.6cm] (a3) {3. 具体任务\\[-1pt]\scriptsize 事实/推断/写作/口述};
\node[cnodegreen,right=of a3,minimum width=2.8cm] (a4) {4. 考试适配器\\[-1pt]\scriptsize 考研/IELTS/TOEFL};

\draw[mainflow] (a1) -- (a2);
\draw[mainflow] (a2) -- (a3);
\draw[altflow] (a3) -- (a4);
\end{tikzpicture}
\end{center}
这里的“考试适配器”指：某一考试从完整语言能力空间中选择哪些任务、采用什么材料、设置什么时间与评分规则。

\subsection{0.2 A1、A2、B1、B2、C1、C2 也不是六套课程}
CEFR 的 A1--C2 是能力水平描述，不是六套固定教材或六种题型。真正变化的是：学习者能够稳定处理的\kw{任务范围、语言复杂度、信息隐含程度、语篇长度、处理速度和独立性}不断扩大。

因此，不应把水平简单理解成“词汇量数字”。更准确的认识是：
\[
\boxed{
\text{Proficiency Level}
=
\text{可稳定完成的语言任务空间}
}
\]

\section{Part I：语言能力的底层结构}

\section{第 1 章：第一层——语言资源是什么？}

\subsection{1.1 语言资源不是“单词 + 语法”}
完整的语言资源至少包含五个由低到高的层级：
\begin{center}
\begin{tikzpicture}[node distance=0.35cm]
\node[cnodeblue,minimum width=2.4cm] (r1) {1. Form\\[-1pt]\scriptsize 声音与拼写};
\node[cnodeorange,right=of r1,minimum width=2.5cm] (r2) {2. Lex/Grammar\\[-1pt]\scriptsize 词汇与句法结构};
\node[cnodeteal,right=of r2,minimum width=2.4cm] (r3) {3. Meaning\\[-1pt]\scriptsize 核心命题与关系};
\node[cnodepurple,right=of r3,minimum width=2.4cm] (r4) {4. Discourse\\[-1pt]\scriptsize 语篇与连贯};
\node[cnodegreen,right=of r4,minimum width=2.5cm] (r5) {5. Pragmatics\\[-1pt]\scriptsize 语用与交际功能};

\draw[mainflow] (r1) -- (r2);
\draw[mainflow] (r2) -- (r3);
\draw[mainflow] (r3) -- (r4);
\draw[altflow] (r4) -- (r5);
\end{tikzpicture}
\end{center}
这些英文术语在本册中分别表示：
\begin{itemize}
  \item \term{Form}{语言形式}：声音、拼写、词形、标点和词序；
  \item \term{Lexicon/Grammar}{词汇与语法}：词义、搭配、配价、句法结构以及组合规则；
  \item \term{Meaning}{意义}：一个短语或句子真正表达的命题内容；
  \item \term{Discourse}{语篇}：多个句子如何通过指代、逻辑、主题与功能组成更大的意义；
  \item \term{Pragmatics}{语用}：谁在什么情境下对谁说什么，以及一句话实际完成什么交际功能。
\end{itemize}

\subsection{1.2 Form：看得见的词和听得见的词是两套入口}
一个单词可能在阅读中认识，在听力中却无法即时识别。原因是书面形式和声音形式虽然指向同一个词汇项目，却通过不同的感知入口进入系统：
\[
\text{Written Form}
\rightarrow
\text{Lexical Entry}
\qquad
\text{Sound Form}
\rightarrow
\text{Lexical Entry}.
\]
因此，“阅读词汇量大但听力识词慢”并不矛盾。真正缺失的是声音形式与词汇意义之间的自动连接，而不是新的中文释义。

\subsection{1.3 Lexicon/Grammar：语言资源必须包含组合规则}
只认识单词不能保证理解句子。阅读或听力必须完成：
\begin{center}
\begin{tikzpicture}[node distance=2.5cm]
\node[cnodeblue,minimum width=3.2cm] (w1) {1. Words\\[-1pt]\scriptsize 词汇与义项};
\node[cnodeorange,right=of w1,minimum width=3.4cm] (w2) {2. Structure\\[-1pt]\scriptsize 主干与句法结构};
\node[cnodegreen,right=of w2,minimum width=3.2cm] (w3) {3. Proposition\\[-1pt]\scriptsize 核心命题表达};

\draw[mainflow] (w1) -- (w2);
\draw[altflow] (w2) -- (w3);
\end{tikzpicture}
\end{center}
例如，能否快速识别主干、从句、修饰范围、否定范围、情态强度，会直接影响后面的意义判断。

\begin{method}[训练诊断]
如果一句长句中的大多数单词都认识，但必须反复回读才能理解，优先怀疑的不是词汇量，而是\kw{结构解析没有自动化}。此时继续背单词的边际收益通常低于针对句法主干、从句边界和修饰关系的训练。
\end{method}

\section{第 2 章：第二层——从句义到语篇}

\subsection{2.1 Meaning：语言处理的最小目标不是翻译，而是恢复命题}
一句话的意义可以先压缩成：
\begin{center}
\begin{tikzpicture}[node distance=2.5cm]
\node[cnodeblue,minimum width=3.2cm] (e1) {1. Entity\\[-1pt]\scriptsize 谁/什么 (对象)};
\node[cnodeorange,right=of e1,minimum width=3.4cm] (e2) {2. Relation\\[-1pt]\scriptsize 建立什么关系 (动作/因果)};
\node[cnodegreen,right=of e2,minimum width=3.2cm] (e3) {3. Qualifier\\[-1pt]\scriptsize 限定范围/时间/程度};

\draw[mainflow] (e1) -- (e2);
\draw[altflow] (e2) -- (e3);
\end{tikzpicture}
\end{center}
即：讨论谁或什么、建立什么关系、有哪些时间/程度/范围/可能性限定。

这比“逐词翻译”更稳定，因为考试中的推断、阅读选项、听力信息点和写作内容都最终依赖命题，而不是中文逐字对应。

\subsection{2.2 Discourse：文章和对话不是独立句子的集合}
语篇至少由四种机制维持：
\begin{center}
\begin{tikzpicture}[node distance=0.40cm]
\node[cnodeblue,minimum width=2.6cm] (d1) {1. Reference\\[-1pt]\scriptsize 指代衔接};
\node[cnodeorange,right=of d1,minimum width=2.6cm] (d2) {2. Relation\\[-1pt]\scriptsize 因果/转折关系};
\node[cnodeteal,right=of d2,minimum width=2.6cm] (d3) {3. Topic\\[-1pt]\scriptsize 主题推进};
\node[cnodegreen,right=of d3,minimum width=2.8cm] (d4) {4. Function\\[-1pt]\scriptsize 语篇功能段落};

\draw[mainflow] (d1) -- (d2);
\draw[mainflow] (d2) -- (d3);
\draw[altflow] (d3) -- (d4);
\end{tikzpicture}
\end{center}
中文分别是指代、句间关系、主题连续和语篇功能。

这意味着一个学习者可以“每句话大概看懂”，却仍然不能准确回答：作者为什么举这个例子、某个代词指什么、某一段与上一段是什么关系。此时缺失的是语篇整合能力，而不是句子翻译能力。

\subsection{2.3 Pragmatics：句子的字面意义不等于实际功能}
\term{语用}{pragmatics} 处理语言与情境、说话者、听话者和交际目的之间的关系。

例如：
\begin{eng}
Could you open the window?
\end{eng}
字面形式是询问能力，实际交际功能通常是请求。写作、口语、态度判断、委婉表达和互动修复都依赖这一层。

\begin{corebox}[语言资源的稳定主线]
无论是 A1 还是 C1，无论考研还是 IELTS，语言资源本身都沿同一主线组织：
\[
\boxed{
\text{形式}
\rightarrow
\text{词汇/语法}
\rightarrow
\text{意义}
\rightarrow
\text{语篇}
\rightarrow
\text{语用}
}
\]
水平提高不是换一套系统，而是在这套系统上扩大范围、提高精确度和自动化程度。
\end{corebox}

\newpage
\section{Part II：在线处理——知道英语为什么仍然用不出来？}

\section{第 3 章：知识与实时执行必须分开}

\subsection{3.1 静态知识不等于在线能力}
一个表达“见过”“背过”“看答案会了”，只说明它存在于某种知识状态中；考试和真实沟通要求的是在有限时间内完成提取和使用。

因此必须区分：
\[
\boxed{
\text{Stored Knowledge}
\neq
\text{Accessible Performance}
}
\]
前者是“我知道”，后者是“任务出现时我能及时调用”。

\subsection{3.2 在线处理链}
语言在真实时间中运行时，可以用下面的过程描述：
\begin{center}
\begin{tikzpicture}[node distance=0.35cm]
\node[cnodeblue,minimum width=2.4cm] (op1) {1. Perceive\\[-1pt]\scriptsize 感知输入};
\node[cnodeorange,right=of op1,minimum width=2.4cm] (op2) {2. Recognize\\[-1pt]\scriptsize 识别形式};
\node[cnodeteal,right=of op2,minimum width=2.4cm] (op3) {3. Construct\\[-1pt]\scriptsize 意义构建};
\node[cnodepurple,right=of op3,minimum width=2.4cm] (op4) {4. Plan/Decide\\[-1pt]\scriptsize 任务计划};
\node[cnodered,right=of op4,minimum width=2.4cm] (op5) {5. Retrieve\\[-1pt]\scriptsize 语言提取};
\node[cnodegreen,right=of op5,minimum width=2.4cm] (op6) {6. Respond\\[-1pt]\scriptsize 做出回答};

\draw[mainflow] (op1) -- (op2);
\draw[mainflow] (op2) -- (op3);
\draw[mainflow] (op3) -- (op4);
\draw[mainflow] (op4) -- (op5);
\draw[altflow] (op5) -- (op6);
\end{tikzpicture}
\end{center}
中文是：感知输入、识别语言单位、构造意义、形成任务计划或判断、提取表达、做出回答。

对阅读选择题，后半段可能是“判断选项”；对口语是“提取表达并说出”；对写作是“形成语篇计划后生成文字”。

\subsection{3.3 自动化为什么重要？}
工作记忆容量有限。如果识别单词、分析结构本身消耗过多注意力，就没有足够资源处理更高层的意义、语篇和任务要求。

因此所谓\term{自动化}{automaticity}，不是“不思考”，而是低层操作足够熟练，从而把注意力释放给高层任务。

\begin{exam}[典型症状与断点]
\begin{center}
\small
\begin{tabularx}{\linewidth}{>{\bfseries\raggedright\arraybackslash}p{0.28\linewidth} Y}
\toprule
表现 & 更可能的断点 \\
\midrule
看到词认识，听到反应不出 & Sound Form $\rightarrow$ Lexical Recognition \\
单词都懂，长句仍反复回读 & Structure Parsing \\
文章大意懂，阅读题仍错 & Meaning $\rightarrow$ Evidence/Task Decision \\
知道中文想法，作文写不出 & Meaning $\rightarrow$ Lexical/Syntactic Retrieval \\
句子能写，整篇作文散乱 & Discourse Planning \\
口语有想法但停顿很长 & Planning/Retrieval $\rightarrow$ Real-time Production \\
听懂前一句，下一句就忘 & 低层识别占用过多 Working Memory \\
\bottomrule
\end{tabularx}
\end{center}
\end{exam}

\section{第 4 章：监控与修复也是语言能力}

\subsection{4.1 完整任务不是一次性直线执行}
成熟语言使用者并不是永远一次成功，而是能够发现沟通是否失败，并采取修复动作：
\begin{center}
\begin{tikzpicture}[node distance=2.5cm]
\node[cnodeblue,minimum width=3.2cm] (pm1) {1. Perform\\[-1pt]\scriptsize 语言任务执行};
\node[cnodeorange,right=of pm1,minimum width=3.4cm] (pm2) {2. Monitor\\[-1pt]\scriptsize 实时效果与风险监控};
\node[cnodegreen,right=of pm2,minimum width=3.2cm] (pm3) {3. Repair\\[-1pt]\scriptsize 改述/澄清/结构重组};

\draw[mainflow] (pm1) -- (pm2);
\draw[altflow] (pm2) -- (pm3);
\end{tikzpicture}
\end{center}
\term{修复}{repair} 指在理解或表达出现问题时，通过重读、改述、澄清、请求重复、降低句式复杂度、重新组织答案等方式恢复任务。

\subsection{4.2 修复不是“能力差的表现”}
在互动中说：
\begin{eng}
Sorry, do you mean ...?\\
Let me put that another way.\\
What I mean is ...
\end{eng}
本质上是控制沟通过程。CEFR 的行动导向框架也把语言策略和沟通活动联系起来，而不是只评价词汇和语法库存。

\newpage
\section{Part III：语言活动——听说读写为什么不应该彼此割裂}

\section{第 5 章：Reception——接收已有语言并恢复意义}

\subsection{5.1 接收包括阅读和听力}
\term{接收}{Reception} 指从外部语言输入中恢复意义。阅读与听力共享后半段处理：
\begin{center}
\begin{tikzpicture}[node distance=0.35cm]
\node[cnodeblue,minimum width=2.4cm] (in1) {Input\\[-1pt]\scriptsize 视/听信号};
\node[cnodeorange,right=of in1,minimum width=2.4cm] (in2) {Form Recognition\\[-1pt]\scriptsize 形式匹配};
\node[cnodeteal,right=of in2,minimum width=2.4cm] (in3) {Structure\\[-1pt]\scriptsize 句法解析};
\node[cnodepurple,right=of in3,minimum width=2.4cm] (in4) {Meaning\\[-1pt]\scriptsize 命题生成};
\node[cnodegreen,right=of in4,minimum width=2.4cm] (in5) {Discourse Model\\[-1pt]\scriptsize 语篇模型};

\draw[mainflow] (in1) -- (in2);
\draw[mainflow] (in2) -- (in3);
\draw[mainflow] (in3) -- (in4);
\draw[altflow] (in4) -- (in5);
\end{tikzpicture}
\end{center}
区别主要在输入入口：阅读面对视觉符号，听力面对连续声流。

\subsection{5.2 阅读的额外优势与听力的额外压力}
阅读允许：停留、回看、重读、利用版面；听力通常具有更强的时间不可逆性，需要持续维护当前语篇状态。

因此听力训练不能简单等于“把阅读材料念出来”。它还需要强化：
\begin{itemize}
  \item 声音边界识别；
  \item 弱读、连读和语流中的词汇识别；
  \item 在声音继续前进时持续更新意义；
  \item 丢失局部信息后快速恢复主线。
\end{itemize}

\section{第 6 章：Production——从意义生成语言}

\subsection{6.1 产出方向与接收方向相反}
\term{产出}{Production} 指学习者主动形成语言。写作和部分独白式口语共享：
\begin{center}
\begin{tikzpicture}[node distance=0.35cm]
\node[cnodeblue,minimum width=2.4cm] (pd1) {Intent\\[-1pt]\scriptsize 意图与目标};
\node[cnodeorange,right=of pd1,minimum width=2.4cm] (pd2) {Meaning\\[-1pt]\scriptsize 命题设计};
\node[cnodeteal,right=of pd2,minimum width=2.4cm] (pd3) {Discourse Plan\\[-1pt]\scriptsize 语篇规划};
\node[cnodepurple,right=of pd3,minimum width=2.4cm] (pd4) {Sentence\\[-1pt]\scriptsize 编码与提取};
\node[cnodegreen,right=of pd4,minimum width=2.4cm] (pd5) {Language Form\\[-1pt]\scriptsize 书面/语音表达};

\draw[mainflow] (pd1) -- (pd2);
\draw[mainflow] (pd2) -- (pd3);
\draw[mainflow] (pd3) -- (pd4);
\draw[altflow] (pd4) -- (pd5);
\end{tikzpicture}
\end{center}

\subsection{6.2 写作与口语的主要差别在时间结构}
写作通常允许更多规划、回看和修改；口语要求实时提取和连续输出。因此两者虽然共享语言资源，却对处理速度和修复方式提出不同要求。

\section{第 7 章：Interaction——互动不是听力加口语}

\subsection{7.1 互动需要维护共同语境}
\term{互动}{Interaction} 是双方不断接收、更新、回应和修复的过程：
\begin{center}
\begin{tikzpicture}[node distance=0.35cm]
\node[cnodeblue,minimum width=2.4cm] (it1) {Listen/Read\\[-1pt]\scriptsize 接收对方};
\node[cnodeorange,right=of it1,minimum width=2.8cm] (it2) {Common Ground\\[-1pt]\scriptsize 更新共同语境};
\node[cnodeteal,right=of it2,minimum width=2.4cm] (it3) {Respond\\[-1pt]\scriptsize 实时回应};
\node[cnodegreen,right=of it3,minimum width=2.6cm] (it4) {Check/Repair\\[-1pt]\scriptsize 检查与修复};

\draw[mainflow] (it1) -- (it2);
\draw[mainflow] (it2) -- (it3);
\draw[altflow] (it3) -- (it4);
\end{tikzpicture}
\end{center}
\term{共同语境}{common ground} 指双方当前共同知道、共同假定和刚刚建立的信息。

\subsection{7.2 互动能力包括调整表达}
同一信息对朋友、教授、陌生人、客户的表达形式不同。互动因此同时调用语言资源、语用判断和实时控制。

\section{第 8 章：Mediation——理解之后重新组织给别人}

\subsection{8.1 中介活动的核心}
CEFR Companion Volume 将\term{中介}{Mediation} 扩展为重要语言活动之一。它不是简单“翻译”，而是：先理解已有信息，再根据新的对象和任务重新选择、组织、解释或转述。\footnote{Council of Europe，\href{https://www.coe.int/en/web/common-european-framework-reference-languages/home}{CEFR Companion Volume 官方页面}。}

其基本过程是：
\begin{center}
\begin{tikzpicture}[node distance=0.40cm]
\node[cnodeblue,minimum width=2.6cm] (me1) {Source\\[-1pt]\scriptsize 来源信息};
\node[cnodeorange,right=of me1,minimum width=2.6cm] (me2) {Meaning Model\\[-1pt]\scriptsize 恢复意义模型};
\node[cnodeteal,right=of me2,minimum width=2.8cm] (me3) {Reorganisation\\[-1pt]\scriptsize 结合新目标重组};
\node[cnodegreen,right=of me3,minimum width=3.0cm] (me4) {Target Output\\[-1pt]\scriptsize 面向目标的再表达};

\draw[mainflow] (me1) -- (me2);
\draw[mainflow] (me2) -- (me3);
\draw[altflow] (me3) -- (me4);
\end{tikzpicture}
\end{center}

\subsection{8.2 常见中介任务}
包括：
\begin{itemize}
  \item 翻译；
  \item 摘要；
  \item paraphrase（改述）；
  \item 根据阅读或听力材料进行综合写作；
  \item 向另一位听众解释复杂信息；
  \item 在不同语言、不同专业背景或不同信息复杂度之间搭桥。
\end{itemize}

\begin{corebox}[语言活动的统一框架]
CEFR Companion Volume 强调 Reception、Production、Interaction、Mediation 等沟通模式，并将学习者视为完成社会行动的语言使用者。\footnote{Council of Europe，\href{https://www.coe.int/en/web/common-european-framework-reference-languages/transparency-and-coherence}{Transparency and coherence}。}
因此，“听说读写”仍然可以作为方便的技能名称，但不应成为最高层的知识架构。
\end{corebox}

\newpage
\section{Part IV：熟练度——A1 到 C2 真正增加了什么？}

\section{第 9 章：能力水平是“稳定工作范围”的扩大}

\subsection{9.1 不要把水平简化成词汇量}
词汇量增加当然重要，但不能解释所有水平差异。更稳定的模型是：
\[
\boxed{
\text{Proficiency}
=
\text{Operating Envelope of Language Use}
}
\]
中文可理解为\kw{语言系统能够稳定工作的任务范围}。

\subsection{9.2 七个扩张轴}
从初级到高级，至少沿以下七个维度扩张：
\begin{enumerate}
  \item \kw{主题范围}：熟悉、具体 $\rightarrow$ 抽象、专业；
  \item \kw{语言范围}：固定表达 $\rightarrow$ 灵活、精确表达；
  \item \kw{输入复杂度}：短、清晰 $\rightarrow$ 长、密集、结构复杂；
  \item \kw{隐含程度}：直接明示 $\rightarrow$ 需要推断与语用理解；
  \item \kw{语篇跨度}：单句 $\rightarrow$ 长篇文本、多轮互动、复杂论证；
  \item \kw{处理压力}：允许缓慢反应 $\rightarrow$ 更自然、连续、实时；
  \item \kw{独立性}：依赖对方帮助 $\rightarrow$ 能自行规划、澄清、修复和调整表达。
\end{enumerate}

\subsection{9.3 能力应该描述成 Profile，而不是单一分数}
一个学习者可能阅读接近 B2、听力 B1、书面产出 B1、实时互动 A2/B1。即使同一技能，在熟悉生活话题和学术任务中也可能表现不同。

因此更准确的表示是：
\[
\boxed{
\text{Proficiency Profile}
=
\text{不同语言活动上的能力轮廓}
}
\]

\begin{method}[诊断意义]
“我的英语是 B1”适合粗略交流；真正设计训练时，应继续拆成：什么任务、什么输入复杂度、什么处理压力、在哪一步第一次失败。
\end{method}

\section{第 10 章：水平提升不是更换策略，而是扩大稳定区}
A1 学习者处理熟悉、直接、短小的任务时，同样调用 Form、Lexicon/Grammar、Meaning 和 Pragmatics；C1 学习者仍然调用同样的结构，只是面对的语言更复杂、意义更隐含、语篇更长、时间压力更高、表达要求更精确。

因此：
\[
\boxed{
\text{Same System}
+
\text{Larger Range}
+
\text{Higher Automaticity}
+
\text{Greater Precision}
}
\]
比“初级方法 / 高级方法完全不同”更适合作为长期学习模型。

\newpage
\section{Part V：考试是怎样从语言能力中取样的？}

\section{第 11 章：考试成绩不是语言能力本身}
一次考试只能通过有限数量的任务观察能力：
\begin{center}
\begin{tikzpicture}[node distance=2.5cm]
\node[cnodeblue,minimum width=3.2cm] (ts1) {1. Proficiency\\[-1pt]\scriptsize 真实语言能力空间};
\node[cnodeorange,right=of ts1,minimum width=3.4cm] (ts2) {2. Performance\\[-1pt]\scriptsize 特定任务/条件下的表现};
\node[cnodegreen,right=of ts2,minimum width=3.2cm] (ts3) {3. Observed Score\\[-1pt]\scriptsize 最终考试得分};

\draw[mainflow] (ts1) -- node[labelnote,above=2pt] {有限抽样} (ts2);
\draw[altflow] (ts2) -- node[labelnote,above=2pt] {按规则量化} (ts3);
\end{tikzpicture}
\end{center}
成绩受到题目抽样、时间限制、任务熟悉度、评分规则和当日执行状态影响。因此，同一人不同考试中的分数不可简单等同。

\subsection{11.1 “考试技巧”真正应该放在哪里？}
稳定语言能力解决“能不能完成这类语言任务”；考试适配器解决：
\begin{itemize}
  \item 当前考试取样哪些任务；
  \item 时间如何分配；
  \item 评分者明确奖励什么；
  \item 题型如何呈现；
  \item 答题格式和交付要求是什么。
\end{itemize}

因此：
\[
\boxed{
\text{Exam Strategy}
=
\text{Language Ability}
+
\text{Task Familiarity}
+
\text{Scoring/Time Adaptation}
}
\]

\section{第 12 章：不同考试只是不同 Adapter}

\subsection{12.1 考研英语}
考研英语主要抽样书面接收、语篇理解、词汇/语法运用、翻译/中介和书面产出。其专题方法可以分别挂接到：
\begin{itemize}
  \item 完形：词项恢复；
  \item 新题型：语篇结构恢复；
  \item 阅读：证据定位与候选命题判断；
  \item 翻译：意义恢复与中文重构；
  \item 写作：任务到语言的生成。
\end{itemize}
它不直接充分测量实时听力与口语互动，因此高分不能自动推出同等强度的口语/听力能力。

\subsection{12.2 IELTS}
IELTS 直接分别测量 Listening、Reading、Writing 和 Speaking。Writing 评分明确包含任务完成/回应、连贯与衔接、词汇资源、语法范围与准确性；Speaking 包含流利与连贯、词汇资源、语法范围与准确性、发音。\footnote{IELTS，\href{https://ielts.org/take-a-test/your-results/ielts-scoring-in-detail}{Understanding your score}。}

从本册模型看，这意味着 IELTS 对\kw{任务完成、语篇组织、语言范围、实时口语控制和发音可理解性}都有直接取样。

\subsection{12.3 TOEFL}
TOEFL 主要面向学术英语沟通。自 2026 年 1 月 21 日起，TOEFL iBT 采用 1--6 分的总分与单项分数尺度，并与 CEFR 更直接对齐；过渡期内仍提供可比的 0--120 总分。\footnote{ETS，\href{https://www.ets.org/toefl/institutions/ibt/score-scale-update.html}{TOEFL iBT Score Scale Update}。}

从统一模型看，TOEFL 的重要意义不是某一具体题型，而是它通过阅读、听力、口语、写作以及综合任务抽样学术场景中的接收、产出和中介能力。题型更新并不会改变这些底层能力的所有权。

\begin{center}
\small
\begin{tabularx}{\linewidth}{>{\bfseries\raggedright\arraybackslash}p{0.18\linewidth} >{\raggedright\arraybackslash}p{0.44\linewidth} Y}
\toprule
体系/考试 & 主要取样重点 & 不应误解为 \\
\midrule
CEFR & 描述不同沟通活动与水平范围 & 一套固定题型 \\
考研英语 & 书面接收、语篇、翻译、书面产出 & 完整英语能力的全部 \\
IELTS & 四技能直接测量，含实时口语 & 只靠题型技巧即可代表能力 \\
TOEFL & 学术场景四技能与综合任务 & 某一固定题型集合永远不变 \\
\bottomrule
\end{tabularx}
\end{center}

\newpage
\section{Part VI：任何英语考试都可以使用的任务控制模型}

\section{第 13 章：先判断“我要完成什么语言行为”}
面对一道题，第一问不应该只是“这是听力还是阅读”，而是：
\[
\boxed{
\text{当前任务要求我对信息做什么操作？}
}
\]
常见操作包括：
\begin{itemize}
  \item Locate：定位事实；
  \item Interpret：解释意义；
  \item Infer：作最小推断；
  \item Compare：比较信息；
  \item Summarise：概括；
  \item Respond：回应；
  \item Explain：解释；
  \item Persuade：说服；
  \item Mediate：转述或重组已有信息。
\end{itemize}

\subsection{13.1 通用五步控制链}
\begin{center}
\begin{tikzpicture}[node distance=0.35cm]
\node[cnodeblue,minimum width=2.4cm] (tk1) {1. Task\\[-1pt]\scriptsize 确认解题目标};
\node[cnodeorange,right=of tk1,minimum width=2.4cm] (tk2) {2. Meaning\\[-1pt]\scriptsize 建立命题模型};
\node[cnodeteal,right=of tk2,minimum width=2.4cm] (tk3) {3. Action\\[-1pt]\scriptsize 定位/推断/概括};
\node[cnodepurple,right=of tk3,minimum width=2.4cm] (tk4) {4. Decision\\[-1pt]\scriptsize 选项/生成表达};
\node[cnodegreen,right=of tk4,minimum width=2.4cm] (tk5) {5. Monitor\\[-1pt]\scriptsize 验证任务完成};

\draw[mainflow] (tk1) -- (tk2);
\draw[mainflow] (tk2) -- (tk3);
\draw[mainflow] (tk3) -- (tk4);
\draw[altflow] (tk4) -- (tk5);
\end{tikzpicture}
\end{center}
中文是：
\begin{enumerate}
  \item \kw{任务}：我要完成什么；
  \item \kw{意义}：输入实际表达什么，或我准备表达什么；
  \item \kw{操作}：我要定位、推断、总结、回应还是组织论证；
  \item \kw{语言/决策}：形成选项判断或生成表达；
  \item \kw{检查}：是否真正完成任务，有无明显证据、语言或交付错误。
\end{enumerate}

\section{第 14 章：接收任务和产出任务方向相反，但控制结构一致}

\subsection{14.1 接收型任务}
\[
\boxed{
\text{Language Input}
\rightarrow
\text{Meaning}
\rightarrow
\text{Task Operation}
\rightarrow
\text{Answer}
}
\]
例如阅读理解、听力选择、词义判断。

\subsection{14.2 产出型任务}
\[
\boxed{
\text{Task/Intent}
\rightarrow
\text{Meaning}
\rightarrow
\text{Discourse Plan}
\rightarrow
\text{Language Output}
}
\]
例如作文、口语回答、邮件。

\subsection{14.3 中介型任务}
\[
\boxed{
\text{Source Input}
\rightarrow
\text{Meaning}
\rightarrow
\text{Selection/Reorganisation}
\rightarrow
\text{New Output}
}
\]
例如翻译、summary、integrated writing。

\begin{exam}[考试中最有价值的元问题]
当一项任务卡住时，不要先问“我英语是不是不行”，而问：
\begin{quote}
\kw{我现在卡在感知、识别、结构、意义、语篇、任务判断、提取表达还是监控修复？}
\end{quote}
这能够直接决定下一步是继续做题、换处理方式，还是在训练后补哪一个部件。
\end{exam}

\newpage
\section{Part VII：统一故障地图——从“英语差”变成可修复断点}

\section{第 15 章：九个高频断点}
\begin{center}
\small
\begin{tabularx}{\linewidth}{>{\bfseries\raggedright\arraybackslash}p{0.22\linewidth} >{\raggedright\arraybackslash}p{0.46\linewidth} Y}
\toprule
断点 & 可观察表现 & 优先训练 \\
\midrule
Signal Recognition & 声音/拼写识别慢 & 高频形式与声音映射 \\
Lexical Access & 认识但提取慢 & 语境化高频检索、主动回忆 \\
Structure Parsing & 长句主干不稳 & 结构拆分与快速组合 \\
Meaning Construction & 句子词都懂但判断错 & 命题压缩、否定/范围/情态 \\
Discourse Integration & 单句懂但文章断裂 & 指代、关系、段落功能 \\
Task Interpretation & 文章懂但答非所问 & 任务限定与目标识别 \\
Planning & 知道观点但组织不出来 & 内容命题与语篇规划 \\
Language Retrieval & 中文想法有，英文出不来 & 功能语块与句型检索 \\
Monitoring / Repair & 错误发生后无法发现/恢复 & 检查清单与修复策略 \\
\bottomrule
\end{tabularx}
\end{center}

\subsection{15.1 第一个断点比最后一个错误更重要}
例如阅读题最终选错，可能并不是“选项判断能力差”，而是更早地误解了否定范围；口语停顿可能并非“没有观点”，而是词汇提取太慢。训练应修复\kw{第一次使后续状态必然错误的节点}。

\section{第 16 章：不要用完整任务掩盖局部故障}

\subsection{16.1 Full Task 负责暴露问题}
完整阅读、听力、口语模拟、作文的价值是让真实问题出现。但如果已经知道断点，继续无限刷完整套题并不会自动修复局部机制。

\subsection{16.2 Isolate Skill 负责修机制}
训练循环应当是：
\begin{center}
\begin{tikzpicture}[node distance=0.35cm]
\node[cnodeblue,minimum width=2.4cm] (tc1) {1. Full Task\\[-1pt]\scriptsize 暴露实际问题};
\node[cnodeorange,right=of tc1,minimum width=2.4cm] (tc2) {2. Breakpoint\\[-1pt]\scriptsize 定位首个断点};
\node[cnodeteal,right=of tc2,minimum width=2.4cm] (tc3) {3. Isolate Skill\\[-1pt]\scriptsize 抽离针对训练};
\node[cnodepurple,right=of tc3,minimum width=2.4cm] (tc4) {4. Practice\\[-1pt]\scriptsize 机制熟练度强化};
\node[cnodegreen,right=of tc4,minimum width=2.4cm] (tc5) {5. Return Task\\[-1pt]\scriptsize 回到完整任务验证};

\draw[mainflow] (tc1) -- (tc2);
\draw[mainflow] (tc2) -- (tc3);
\draw[mainflow] (tc3) -- (tc4);
\draw[altflow] (tc4) -- (tc5);
\end{tikzpicture}
\end{center}
例如：
\begin{itemize}
  \item 听力错很多，诊断为弱读下的识词失败：先练 Sound $\rightarrow$ Word，再回完整听力；
  \item 阅读长句反复误判：先练 Structure $\rightarrow$ Proposition，再回整篇；
  \item 写作句子都对但段落散：先练 Proposition $\rightarrow$ Discourse Plan，再回完整作文。
\end{itemize}

\begin{warn}[专项训练的边界]
专项训练只负责修复局部机制，不能永久替代完整任务。若长期只背词、拆句、做跟读或练模板，而不重新回到真实语言任务，技能之间不会自动重新整合。
\end{warn}

\newpage
\section{Part VIII：怎样设计个人英语学习系统}

\section{第 17 章：先建立能力 Profile，再制定计划}

\subsection{17.1 诊断不要只记录分数}
每个主要任务至少记录四类信息：
\begin{enumerate}
  \item \kw{Task}：做的是什么任务；
  \item \kw{Outcome}：正确率、完成度、流利度或评分结果；
  \item \kw{First Breakpoint}：第一个控制失败点；
  \item \kw{Repair}：下一阶段的具体训练动作。
\end{enumerate}

\subsection{17.2 学习计划应该围绕“瓶颈”而不是平均分配}
如果阅读已经稳定，而听力的 Sound $\rightarrow$ Word 仍然很弱，继续平均分配四项时间并不高效。训练资源应优先投向当前最限制目标考试或真实沟通任务的瓶颈。

\subsection{17.3 目标必须写成可执行任务}
比“我要提高英语到 B2”更有训练价值的是：
\begin{itemize}
  \item 能听懂正常语速下熟悉学术话题的主线和主要论据；
  \item 能在 90 秒内就熟悉话题连续表达并自行修复明显停顿；
  \item 能读完一篇中等复杂度文章并准确恢复段落关系；
  \item 能写一封完成全部交际动作、语气得体的正式邮件。
\end{itemize}
任务一旦具体，断点才可观察，训练才可设计。

\section{第 18 章：输入、输出和反馈怎样形成闭环}

\subsection{18.1 输入不是“接触英语越多越好”}
高价值输入应该兼顾：
\begin{itemize}
  \item 当前大部分可理解；
  \item 含有略高于当前稳定区的结构或表达；
  \item 能被重复处理；
  \item 能与当前目标任务发生迁移。
\end{itemize}

\subsection{18.2 输出的价值是暴露检索与组织缺口}
只有阅读词汇表，很难知道哪些词真正能用于产出。写作和口语会暴露：
\begin{itemize}
  \item 词汇提取不出；
  \item 搭配不自然；
  \item 句型过度单一；
  \item 语篇没有计划；
  \item 语体不合适。
\end{itemize}

\subsection{18.3 反馈应该指向机制}
低价值反馈：
\begin{quote}
“这道题错了，答案是 B。”
\end{quote}
高价值反馈：
\begin{quote}
“你把 may 读成 must，导致 Meaning Construction 阶段把可能性错误升级为必然性；后面的选项判断只是结果。”
\end{quote}

\newpage
\section{Part IX：现有专题手册如何挂接到总图}

\section{第 19 章：英语专题的 Canonical Ownership}
\begin{center}
\small
\begin{tabularx}{\linewidth}{>{\bfseries\raggedright\arraybackslash}p{0.18\linewidth} >{\raggedright\arraybackslash}p{0.46\linewidth} Y}
\toprule
专题 & 主要拥有的机制 & 与总图的接口 \\
\midrule
发音/听音 & Sound Form 与识别 & Form $\rightarrow$ Lexical Access \\
词汇/搭配 & Lexicon、Collocation、Valency & Form/Meaning 资源 \\
语法/长难句 & Structure Parsing & Lexicon/Grammar $\rightarrow$ Meaning \\
完形 & 多层约束下恢复词项 & Context $\rightarrow$ Missing Form \\
新题型 & 衔接、连贯、语篇结构 & Discourse Relations \\
阅读 & 证据定位与候选命题判断 & Meaning $\rightarrow$ Task Decision \\
翻译 & 意义恢复与重新编码 & Reception $\rightarrow$ Mediation \\
写作 & 从任务到书面语言生成 & Meaning $\rightarrow$ Production \\
口语 & 实时计划与输出 & Meaning $\rightarrow$ Real-time Production \\
互动 & 共同语境与修复 & Reception $\leftrightarrow$ Production \\
\bottomrule
\end{tabularx}
\end{center}

\subsection{19.1 总图不重复专题正文}
例如总图只说明“阅读题需要 Meaning $\rightarrow$ Evidence $\rightarrow$ Decision”，具体怎样定位证据、比较选项，应由阅读手册拥有；总图只说明写作属于 Production，具体如何构造命题和段落，应由写作手册拥有。

这保证每个概念只有一个主要 Owner，同时保持专题可以通过接口互相调用。

\section{第 20 章：跨技能迁移应该明确说明“迁移什么”}
不要泛泛地说“阅读好会带动写作”。真正可能迁移的是具体机制：
\begin{itemize}
  \item 阅读中的句法解析可以支持听力中的句义构造；
  \item 新题型中的语篇关系识别可以支持阅读主旨和写作连贯；
  \item 写作中的命题组织可以支持口语长回答；
  \item 听力中的声音识别不会因为阅读词汇量大而自动形成，必须额外训练声形映射。
\end{itemize}

\newpage
\section{Part X：最终压缩页}

\section{第 21 章：一张图理解英语能力}

\begin{corebox}[语言资源]
\[
\boxed{
\text{Form}
\rightarrow
\text{Lexicon/Grammar}
\rightarrow
\text{Meaning}
\rightarrow
\text{Discourse}
\rightarrow
\text{Pragmatics}
}
\]
解决：\kw{你拥有什么语言资源。}
\end{corebox}

\begin{statebox}[语言活动]
\[
\boxed{
\text{Reception}
\quad
\text{Production}
\quad
\text{Interaction}
\quad
\text{Mediation}
}
\]
解决：\kw{你正在用语言完成什么类型的活动。}
\end{statebox}

\begin{exam}[在线控制]
\[
\boxed{
\text{Perceive}
\rightarrow
\text{Recognize}
\rightarrow
\text{Understand}
\rightarrow
\text{Plan/Decide}
\rightarrow
\text{Respond}
\rightarrow
\text{Repair}
}
\]
解决：\kw{这些资源在真实时间里能否顺利运行。}
\end{exam}

\begin{bridge}[考试控制]
\[
\boxed{
\text{Task}
\rightarrow
\text{Meaning}
\rightarrow
\text{Action}
\rightarrow
\text{Language/Decision}
\rightarrow
\text{Monitor}
}
\]
解决：\kw{当前考试任务下一步具体做什么。}
\end{bridge}

\section{第 22 章：学习者每天真正需要问的五个问题}
\begin{enumerate}
  \item \kw{我当前的目标任务是什么？}不是笼统“学英语”，而是明确到阅读、听懂、对话、写邮件、做综合题等；
  \item \kw{完成这个任务最先失败在哪？}声音、词汇、句法、意义、语篇、任务理解、提取还是监控；
  \item \kw{这个断点需要哪一种最小专项训练？}只训练当前机制；
  \item \kw{什么时候回到完整任务验证？}专项必须重新整合；
  \item \kw{新任务中是否仍出现同一断点？}只有跨样本稳定改善，能力才真正更新。
\end{enumerate}

\begin{corebox}[最终心智模型]
英语学习不是收集技巧，而是在同一套语言系统上不断扩大\kw{可稳定完成的任务空间}。

真正的长期循环是：
\[
\boxed{
\text{做真实任务}
\rightarrow
\text{找到第一个断点}
\rightarrow
\text{隔离并训练机制}
\rightarrow
\text{回到真实任务}
\rightarrow
\text{扩大稳定工作范围}
}
\]
考试变化时，只更换任务与评分适配器；语言资源、在线处理、沟通活动和诊断逻辑保持稳定。
\end{corebox}

\section{附录 A：权威框架与当前考试信息说明}
本册的能力结构主要参考 CEFR 及其 Companion Volume 对语言活动、策略和沟通能力的组织方式；具体考试示例只用于说明“同一能力系统如何被不同考试取样”，不作为本册母模型的来源。

\begin{itemize}
  \item \href{https://www.coe.int/en/web/common-european-framework-reference-languages/home}{Council of Europe: CEFR Companion Volume 官方页面}
  \item \href{https://www.coe.int/en/web/common-european-framework-reference-languages/transparency-and-coherence}{Council of Europe: Transparency and coherence}
  \item \href{https://ielts.org/take-a-test/your-results/ielts-scoring-in-detail}{IELTS: Understanding your score}
  \item \href{https://www.ets.org/toefl/institutions/ibt/score-scale-update.html}{ETS: TOEFL iBT Score Scale Update（2026-01-21 起）}
\end{itemize}

\section{附录 B：术语速查}
\small
\begin{longtable}{>{\bfseries}p{0.23\linewidth} p{0.25\linewidth} p{0.44\linewidth}}
\toprule
术语 & 中文 & 在本册中的作用 \\
\midrule
Form & 语言形式 & 声音、拼写、词形等可感知入口 \\
Lexicon/Grammar & 词汇与语法 & 词义、搭配、配价和组合结构 \\
Meaning & 意义 & 恢复或生成命题内容 \\
Discourse & 语篇 & 多句之间的指代、关系、主题和功能 \\
Pragmatics & 语用 & 情境、说话者、听者与交际功能 \\
Reception & 接收 & 从阅读/听力输入恢复意义 \\
Production & 产出 & 从意图和意义生成语言 \\
Interaction & 互动 & 双方实时更新共同语境并回应 \\
Mediation & 中介 & 理解来源信息后重组为新的目标输出 \\
Automaticity & 自动化 & 低层处理速度足够快，释放高层注意力 \\
Repair & 修复 & 发现沟通失败后调整理解或表达 \\
Proficiency Profile & 能力轮廓 & 不同语言任务上的能力分布 \\
Exam Adapter & 考试适配器 & 某考试的任务、时间、材料与评分方式 \\
\bottomrule
\end{longtable}
\normalsize

\end{document}
